针对新疆光伏场站复杂天气下功率快速变化及预警可靠性不足的问题，研究多源气象驱动的多尺度概率预测与爬坡风险预警。以预测发布时刻实际可获得的信息为约束，比较不同气象源、时间尺度和预测时距的增益；建立多尺度分位数模型，在真实标签到达后开展滚动校准，评价天气转换中的覆盖偏差与区间效率；保持逐时距边际分布一致，通过独立排序、总体历史秩和天气条件历史秩构造轨迹，检验时间依赖对爬坡事件概率及告警效果的贡献。采用时间独立的数据划分、同预算消融和事件级统计，形成数据处理与评价流程、概率预测实现及离线预警原型。研究为新疆光伏气象服务中的方法选择和风险识别提供可复核依据；地域化结论以取得授权的新疆实测数据为前提，数据与合作条件尚待落实。
